\documentclass{article}

\textwidth=6.5in
\textheight=8.5in
\voffset=-54pt
\oddsidemargin=0in
\evensidemargin=0in
\setlength{\parskip}{8pt}
\setlength{\parindent}{0pt}

\usepackage{workbook}

\usepackage[sols]{handout}

\begin{document}

\begin{center}
  \large\textsc{Problem Set 31 - Slicing to Find Area}
  
      \pspace{12pt}
  \begin{problemonly}\large Name: \rule{5cm}{0.15mm}\\ \end{problemonly}
\end{center}

\begin{grading}
  Total: 25 points
  
    \begin{enumerate}
    \item 20 points
    \item 9 points
    \item 8 points
    \item 4 points
    \end{enumerate}
    Total: 41 points
  
\end{grading}

\begin{enumerate}


\item
\begin{grading}
    20 points: 2/2/2/1/1/2/2/1/1/6
\end{grading}

For help with these problems, read pp. 843-850 in Gottlieb.

  Let $\mathcal{R}$ be the region enclosed by the curves $y = \sqrt{x}$
  and $y = \dfrac{x}{2}$.  In this problem, you'll find the area of
  $\mathcal{R}$ two ways, once by slicing vertically and once by slicing
  horizontally.

  \begin{enumerate}
  \item
    \begin{problem}[\stretch{3}]
      Sketch the region.
    \end{problem}

    \begin{solution}
      The curves $y = \sqrt{x}$ and $y = \frac{x}{2}$ intersect when
      \begin{align*}
        \sqrt{x} &= \frac{x}{2} \\
        2 \sqrt{x} &= x \\
        \shortintertext{Squaring both sides,}
        4 x &= x^2 \\
        4x - x^2 &= 0 \\
        x (4 - x) &= 0 \\
        x &= 0, 4
      \end{align*}
      So, here's a sketch of the region: 
      \begin{center}
        \begin{tikzpicture}
          \begin{axis}[width=3in, height=2in, xtick={4}, ytick={2}]
            \addplot+[domain=0:4, fillregion, draw=none] {sqrt(x)} -- (0,
            0);
            \addplot+[domain=0:4.5] {sqrt(x)};
            \addplot+[domain=0:4.5] {x/2};
          \end{axis}
        \end{tikzpicture}
      \end{center}
    \end{solution}

  \end{enumerate}
  We'll first tackle slicing vertically.  
  \begin{enumerate}[resume]

  \item
    \begin{problem}[\vfill]
      If you slice the region $\mathcal{R}$ vertically into $n$ slices of
      equal thickness $\Delta x$ and label $x_0, x_1, \ldots, x_n$ as
      usual, what is $\Delta x$ in terms of $n$?  What is $x_k$ in terms
      of $\Delta x$?
    \end{problem}

    \begin{solution}
      We're slicing the interval $[0, 4]$, so $\Delta x = \frac{4 - 0}{n}
      = \frac{4}{n}$ and $x_k = 0 + k \Delta x$.
    \end{solution}

  \item

    \begin{problem}[\vfill]
      On your sketch of the region, draw a rectangle that approximates the
      $k$-th slice.  What is the area of this rectangle?
    \end{problem}

    \begin{solution}
      \begin{center}
        \begin{tikzpicture}[declare function={f(\x)=sqrt(\x); g(\x) =
            \x/2;}] 
          \begin{axis}[width=3in, height=2in, xtick=\empty, ytick=\empty,
            axis on top=true, ymin=0, axis x line=bottom, xtick={0,
              0.307692, ..., 4.01}, xticklabels={$x_0$, $x_1$,
              $\cdots$,,,,,$x_k$,,,,,,$x_n$}, disabledatascaling]
            \addplot+[domain=0:4, fillregion, draw=none] {sqrt(x)} -- (0,
            0); %
            \addplot+[vertslices] coordinates { (0, {f(0)}) (4/13,
              {f(4/13)}) (8/13, {f(8/13)}) (12/13, {f(12/13)}) (16/13,
              {f(16/13)}) (20/13, {f(20/13)}) (24/13, {f(24/13)}) (28/13,
              {f(28/13)}) (32/13, {f(32/13)}) (36/13, {f(36/13)}) (40/13,
              {f(40/13)}) (44/13, {f(44/13)}) (48/13, {f(48/13)}) }; %
            \addplot+[fill=white, draw=none, domain=0:4] {x/2}
            \closedcycle; %
            \addplot+[domain=0:4.5] {sqrt(x)}; %
            \addplot+[domain=0:4.5] {x/2}; %
            \draw[fill=cyan, fill opacity=0.75, draw=blue] (24/13,
            {f(28/13)}) -- (28/13, {f(28/13)}) -- (28/13, {g(28/13)}) --
            (24/13, {g(28/13)}) -- cycle; %
          \end{axis}
        \end{tikzpicture}
      \end{center}
      This rectangle's width is $\Delta x$, and its height is (top
      $y$-coordinate) $-$ (bottom $y$-coordinate) = $\sqrt{x_k} -
      \frac{x_k}{2}$.  Therefore, the rectangle's area is $\boxed{\left(
          \sqrt{x_k} - \frac{x_k}{2} \right) \Delta x}$.
    \end{solution}

  \item
    \begin{problem}[\vfill]
      Write a Riemann sum that approximates the area of $\mathcal{R}$.
    \end{problem}

    \begin{solution}
      To approximate the area of the whole region $\mathcal{R}$, we just
      add up our approximations for the areas of all the slices, giving
      $\boxed{\sum_{k=1}^n \left( \sqrt{x_k} - \frac{x_k}{2} \right)
        \Delta x}$.
    \end{solution}

  \item
\begin{problem}[\vfill\newpage]
      Write a definite integral that gives the exact area of
      $\mathcal{R}$. 
    \end{problem}

    \begin{solution}
      Taking the limit as $n \to \infty$ of our Riemann sum, the exact
      area of $\mathcal{R}$ is $\boxed{\int_0^4 \left( \sqrt{x} -
          \frac{x}{2} \right) dx}$. 
    \end{solution}

  \end{enumerate}

  Now, let's instead slice horizontally.
  \begin{enumerate}[resume]
  \item
    \begin{problem}[\vfill]
      If you slice the region $\mathcal{R}$ horizontally into $n$ slices
      of equal thickness $\Delta y$ and label $y_0, y_1, \ldots, y_n$ as
      usual, what is $\Delta y$ in terms of $n$?  What is $y_k$ in terms
      of $\Delta y$?
    \end{problem}

    \begin{solution}
      Now, we're slicing $[0, 2]$ on the $y$-axis, so $\Delta y = \frac{2
        - 0}{n} = \frac{2}{n}$ and $y_k = 0 + k \Delta y$.
    \end{solution}

  \item
    \begin{problem}[\vfill]
      On your sketch of the region, draw a rectangle that approximates the
      $k$-th slice.  What is the area of this rectangle?
    \end{problem}

    \begin{solution}
      \begin{center}
        \begin{tikzpicture}[declare function={f(\x)=2*\x; g(\x)=\x^2;}] 
          \begin{axis}[width=3in, height=2in, xtick=\empty, ytick=\empty,
            axis on top=true, ymin=0, axis x line=bottom, ytick={0,
              0.285714, ..., 2.01}, yticklabels={$y_0 = 0$,,,,,$y_k$,,$y_n
              = 2$}, disabledatascaling]
            \addplot+[domain=0:4, fillregion, draw=none] {sqrt(x)} -- (0,
            0); %
            \addplot+[horizslices] coordinates { ({f(0)}, 0) ({f(2/7)},
              2/7) ({f(4/7)}, 4/7) ({f(6/7)}, 6/7) ({f(8/7)}, 8/7)
              ({f(10/7)}, 10/7) ({f(12/7)}, 12/7) ({f(2)}, 2) };
            \addplot+[fill=white, draw=none, domain=0:4] {sqrt(x)} -- (0,
            2) -- (0, 0); %
            \addplot+[domain=0:4.5] {sqrt(x)}; %
            \addplot+[domain=0:4.5] {x/2}; %
            \draw[fill=cyan, fill opacity=0.75, draw=blue] ({f(10/7)},
            10/7) -- ({g(10/7)}, 10/7) -- ({g(10/7)}, 8/7) -- ({f(10/7)},
            8/7) -- cycle; %
          \end{axis}
        \end{tikzpicture}
      \end{center}
      The height of this slice is $\Delta y$, and its width is (right
      $x$-coordinate) $-$ (left $x$-coordinate).  The left end of the
      slice is on the curve $y = \sqrt{x}$ (or $x = y^2$), and the right
      end is on the curve $y = \frac{x}{2}$ (or $x = 2y$), so the area of
      the rectangle is $\boxed{\left(2 y_k - y_k^2 \right) \Delta y}$.
    \end{solution}

  \item
    \begin{problem}[\vfill]
      Write a Riemann sum that approximates the area of $\mathcal{R}$.
    \end{problem}

    \begin{solution}
      Summing the approximate areas of all the slices, the area of
      $\mathcal{R}$ is approximately $\boxed{\sum_{k=1}^n \left(2 y_k -
          y_k^2 \right) \Delta y}$.
    \end{solution}

  \item
\begin{problem}[\vfill]
      Write a definite integral that gives the exact area of
      $\mathcal{R}$. 
    \end{problem}

    \begin{solution}
      Taking the limit as $n \to \infty$ of our Riemann sum, the exact
      area of $\mathcal{R}$ is $\boxed{ \int_0^2 \left( 2y - y^2 \right)
        dy}$. 
    \end{solution}

  \item
    \begin{problem}[\nstretch{3}]
      Evaluate the integrals you got in
      {{{{\#1}}(e)}} and
      {{{{\#1}}(i)}} to
      show that you get the same answer either way.
    \end{problem}

    \begin{solution}
      Our integral in
      {{{{\#1}}(e)}} was 
      \begin{align*}
        {\int_0^4 \left( \sqrt{x} - \frac{x}{2} \right) dx} &=
        \left. \frac{2}{3} x^{3/2} - \frac{x^2}{4} \right|_0^4 \\
        &= \left( \frac{2}{3} \cdot 4^{3/2} - 4 \right) - 0 \\
        &= \frac{2}{3} \cdot 8 - 4 \\
        &= \boxed{\frac{4}{3}}
      \end{align*}
      Our integral in
      {{{{\#1}}(i)}} was
      \begin{align*}
        { \int_0^2 \left( 2y - y^2 \right) dy} &= \left. y^2 -
          \frac{1}{3}y^3 \right|_0^2 \\
        &= \left( 4 - \frac{8}{3} \right) - 0 \\
        &= \frac{4}{3}
      \end{align*}      
    \end{solution}

  \end{enumerate}

\item

\begin{grading}
    8 points - consult Gradescope rubric
\end{grading}

\begin{problem}[\vfill\newpage]
    Let $\mathcal{R}$ be the region in the first quadrant to the right of
    $y = x^2$, the left of $x + y = 6$, and above $y = 1$.  Find the area
    of $\mathcal{R}$; it's up to you whether to slice vertically or
    horizontally.  Please include a sketch of the region $\mathcal{R}$.

    (Do evaluate your integral to get a numerical answer.)
  \end{problem}

  \begin{solution}
    Here's a picture of the region:
    \begin{center}
      \begin{tikzpicture}
        \begin{axis}[ymax=6, disabledatascaling, axis equal=true,
          width=2.5in, height=2.5in, enlargelimits=true, xtick={6},
          ytick={1, 6}, cycle list name=mycolor, clip=false, cycle list
          shift=2]  
          \addplot+[domain=1:2, fillregion, draw=none]{x^2} -- (5, 1) --
          (1, 1); %
          \addplot+[domain=0:sqrt(6)]{x^2} node[above]{$y=x^2$}; %
          \addplot+[domain=0:6]{6-x} node[pos=0.5, anchor=south west]{$x +
            y = 6$}; %
          \addplot+[domain=0:6]{1} node[right]{$y = 1$}; %
        \end{axis}
      \end{tikzpicture}
    \end{center}
    The top point of the region is where $x^2 = 6 - x$.  This equation can
    be rearranged as $x^2 + x - 6 = 0$ and factored as $(x + 3)(x - 2) =
    0$, and the relevant solution is $x = 2$.  So, the top point is $(2,
    4)$.  The bottom left corner of the region is the point $(1, 1)$, and
    the bottom right is $(5, 1)$.

    If we slice vertically, there are two ``types'' of slices.  Slices to
    the left of $x = 2$ have their top on the curve $y = x^2$, while
    slices to the right have their top on the line $x + y = 6$.  On the
    other hand, if we slice horizontally, then there is only one ``type''
    of slice: all slices have their left side on $y = x^2$ and their right
    side on $x + y = 6$.  So, let's slice horizontally; as usual, we focus
    on the $k$-th slice and approximate it by a rectangle:
    \begin{center}
      \begin{tikzpicture}[declare function={f(\x)=6-\x;}]
        \begin{axis}[ymax=6, disabledatascaling, axis equal=true,
          width=2.5in, height=2.5in, enlargelimits=true, xtick={6},
          ytick={1, 1.333, ..., 4, 6}, yticklabels={$y_0 =
            1$,,,,,$y_k$,,,,$y_n = 4$, $6$}, cycle list name=mycolor,
          clip=false, axis on top=true]
          \addplot+[domain=1:2, fillregion, draw=none]{x^2} -- (5, 1) --
          (1, 1); %
          \addplot+[horizslices, black] coordinates { ({f(1)}, 1)
            ({f(4/3)}, 4/3) ({f(5/3)}, 5/3) ({f(2)}, 2) ({f(7/3)}, 7/3)
            ({f(8/3)}, 8/3) ({f(3)}, 3) ({f(10/3)}, 10/3) ({f(11/3)},
            11/3)}; %
          \addplot+[domain=1:2, fill=white, draw=none] {x^2} -- (0, 4) --
          (0, 1) -- (1, 1); % 
          \addplot+[domain=0:sqrt(6)]{x^2} node[above]{$y=x^2$}; %
          \addplot+[domain=0:6]{6-x} node[pos=0.5, anchor=south west]{$x +
            y = 6$}; %
          \addplot+[domain=0:6]{1} node[right]{$y = 1$}; %
          \draw[blue, fill=cyan, fill opacity=0.75] ({sqrt(8/3)}, 8/3) --
          ({sqrt(8/3)}, 7/3) -- (10/3, 7/3) -- (10/3, 8/3) -- cycle; % 
        \end{axis}        
      \end{tikzpicture}
    \end{center}
    The rectangle has height $\Delta y$ and width (right $x$-coordinate)
    $-$ (left $x$-coordinate).  The right end is on the line $x + y = 6$
    (which can be rewritten as $x = 6 - y$), and the left end is on the
    curve $y = x^2$ (which can be rewritten as $x = \sqrt{y}$).  So,
    \[ \text{area of $k$-th slice} \approx \left( 6 - y_k - \sqrt{y_k}
    \right) \Delta y.\] Summing over all slices and taking the limit as
    the number of slices $\to \infty$, the area of $\mathcal{R}$ is
    \begin{align*}
      \int_1^4 \left( 6 - y - \sqrt{y} \right) dy &= \left. 6y -
        \frac{y^2}{2} - \frac{2}{3} y^{3/2} \right|_1^4 \\
      &= \left( 24 - 8 - \frac{2}{3} \cdot 8 \right) - \left( 6 -
        \frac{1}{2} - \frac{2}{3} \right) \\
      &= \boxed{\frac{35}{6}}
    \end{align*}
  \end{solution}



% \item

% \begin{problem}[\vfill\newpage]
%     In this problem, you'll evaluate $\displaystyle
%     \int_0^{\sqrt{3}/2} \arcsin x\,dx$ by
%     visualizing it as an area and slicing this area horizontally to get an
%      integral which is easier to evaluate than integrating by parts.
%   \end{problem}

%   \begin{solution}
%     We can visualize $\displaystyle \int_0^{\sqrt{3}/2} \arcsin x\,dx$ as
%     the area under $y = \arcsin x$ from $x = 0$ to $x =
%     \frac{\sqrt{3}}{2}$:
%     \begin{center}
%       \begin{tikzpicture}
%         \begin{axis}[xtick={0.866}, xticklabels={$\frac{\sqrt{3}}{2}$},
%           ytick={1.047}, yticklabels={$\frac{\pi}{3}$}, clip=false,
%           width=2in, height=2.5in]
%           \addplot+[domain=0:sqrt(3)/2, fillregion] {rad(asin(x))}
%           \closedcycle; %
%           \addplot+[domain=-pi/2:pi/2] ({sin(deg(x))}, x)
%           node[above] {$y = \arcsin x$};
%         \end{axis}
%       \end{tikzpicture}
%     \end{center}
%     (We've labeled $\arcsin \left(\frac{\sqrt{3}}{2} \right) =
%     \frac{\pi}{3}$ on the $y$-axis, as that's the highest point in the
%     region.)  We're going to slice horizontally, which means we're slicing
%     the interval $\left[ 0, \frac{\pi}{3} \right]$ on the $y$-axis into
%     $n$ slices of equal height $\Delta y$.  We'll label $y_0 = 0, \ldots,
%     y_n = \frac{\pi}{3}$ as usual and focus on the $k$-th slice, which we
%     approximate by a rectangle:
%     \begin{center}
%       \begin{tikzpicture}
%         \begin{axis}[xtick={0.866}, xticklabels={$\frac{\sqrt{3}}{2}$},
%           ytick={1.047, 0.698}, yticklabels={$\frac{\pi}{3}$, $y_k$}, axis
%           on top, disabledatascaling, clip=false, width=2in, height=2.5in]
%           \addplot+[domain=0:sqrt(3)/2, fillregion] {rad(asin(x))} \closedcycle;
%           \addplot[horizslices] coordinates {
%             ({sqrt(3)/2}, 0) ({sqrt(3)/2}, pi/27) ({sqrt(3)/2}, 2*pi/27)
%             ({sqrt(3)/2}, 3*pi/27) ({sqrt(3)/2}, 4*pi/27) ({sqrt(3)/2}, 5*pi/27)
%             ({sqrt(3)/2}, 6*pi/27) ({sqrt(3)/2}, 7*pi/27) ({sqrt(3)/2}, 8*pi/27)
%             ({sqrt(3)/2}, 9*pi/27)};
%           \addplot+[domain=0:sqrt(3)/2, fill=white, draw=none]
%           {rad(asin(x))} -- (0, {pi/3}) -- (0, 0);
%           \addplot+[domain=-pi/2:pi/2] ({sin(deg(x))}, x)
%           node[above] {$y = \arcsin x$};
%           \draw[fill=cyan, draw=blue, fill opacity=0.75] ({sqrt(3)/2},
%           5*pi/27) -- ({sin(deg(6*pi/27))}, 5*pi/27) --
%           ({sin(deg(6*pi/27))}, 6*pi/27) -- ({sqrt(3)/2},
%           6*pi/27) -- cycle;
%           \draw[dashed] (0, 6*pi/27) -- ({sin(deg(6*pi/27))}, 6*pi/27);
%         \end{axis}
%       \end{tikzpicture}
%     \end{center}
%     This rectangle's height is $\Delta y$, and its width is (right
%     $x$-value) $-$ (left $x$-value).  The right $x$-value is
%     $\frac{\sqrt{3}}{2}$; since $y = \arcsin x$ can be rewritten as $x =
%     \sin y$, the left $x$-value is $\sin (y_k)$.  Therefore,
%     \[ \text{area of $k$-th slice} \approx \left( \frac{\sqrt{3}}{2} -
%       \sin y_k \right) \Delta y.\] Summing the areas of all the slices and
%     taking the limit as the number of slices $\to \infty$, the area of the
%     region is 
%     \begin{align*}
%       \int_0^{\pi/3} \left( \frac{\sqrt{3}}{2} - \sin y \right) dy &=
%       \left. \frac{\sqrt{3}}{2} y + \cos y \right|_0^{\pi/3} \\
%       &= \left( \frac{\sqrt{3}}{2} \cdot \frac{\pi}{3} + \frac{1}{2}
%       \right) - (0 + 1) \\
%       &= \boxed{\frac{\pi \sqrt{3}}{6} - \frac{1}{2}}
%     \end{align*}
%   \end{solution}

\item

\begin{grading}
    9 points - consult Gradescope rubric
\end{grading}
  \begin{problem}[\vfill\newpage]
    In this problem, evaluate $\displaystyle
    \int_0^{1} \arccos x\,dx$ by first
    visualizing it as an area and then slicing this area horizontally to get an
     integral which is easier to evaluate than integrating by parts. Then, evaluate the resulting integral.
  \end{problem}

  \begin{solution}
    We can visualize $\displaystyle \int_0^{1} \arccos x\,dx$ as
    the area under $y = \arccos x$ from $x = 0$ to $x =
    1$:
    \begin{center}
      \begin{tikzpicture}
        \begin{axis}[xtick={1}, xticklabels={$1$},
          ytick={1.571}, yticklabels={$\frac{\pi}{2}$}, clip=false,
          width=2in, height=2.5in]
          \addplot+[domain=0:1, fillregion] {rad(acos(x))}
          \closedcycle; %
          \addplot+[domain=0:pi] ({cos(deg(x))}, x)
          node[above] {$y = \arccos x$};
        \end{axis}
      \end{tikzpicture}
    \end{center}
    (We've labeled $\arccos \left(0\right) =
    \frac{\pi}{2}$ on the $y$-axis, as that's the highest point in the
    region.)  We're going to slice horizontally, which means we're slicing
    the interval $\left[ 0, \frac{\pi}{2} \right]$ on the $y$-axis into
    $n$ slices of equal height $\Delta y$.  We'll label $y_0 = 0, \ldots,
    y_n = \frac{\pi}{2}$ as usual and focus on the $k$-th slice, which we
    approximate by a rectangle:
    \begin{center}
      \begin{tikzpicture}
        \begin{axis}[xtick={1},
          ytick={0.589, 1.571}, yticklabels={$y_k$,$\frac{\pi}{2}$}, axis
          on top, disabledatascaling, clip=false, width=2in, height=2.5in]
          \addplot+[domain=0:1, fillregion] {rad(acos(x))} \closedcycle;
          \addplot[horizslices] coordinates {
            (1, 0) (1, pi/16) (1, 2*pi/16)
            (1, 3*pi/16) (1, 4*pi/16) (1, 5*pi/16)
            (1, 6*pi/16) (1, 7*pi/16) (1, 8*pi/16)};
          \addplot+[domain=0:1, fill=white, draw=none]
          {rad(acos(x))} -- (1.1, {pi/2+0.1}) -- (0, {pi/2+0.1});
          \addplot+[domain=0:pi/2, very thick] ({cos(deg(x))}, x)
          node[midway, right] {$x = \cos y$};
          \draw[fill=cyan, draw=blue, fill opacity=0.75] (0,
          2*pi/16) -- ({cos(deg(3*pi/16))}, 2*pi/16) --
          ({cos(deg(3*pi/16))}, 3*pi/16) -- (0,
          3*pi/16) -- cycle;
          \draw[dashed] (0, 3*pi/16) -- ({cos(deg(3*pi/16))}, 3*pi/16);
        \end{axis}
      \end{tikzpicture}
    \end{center}
    This rectangle's height is $\Delta y$, and its width is (right
    $x$-value) $-$ (left $x$-value).  The right $x$-value is $\cos(y_k)$, since $y=\arccos(x)$ can be rewritten as $x=\cos(y)$ and the left $x$-value is 0. 
  Therefore,
    \[ \text{area of $k$-th slice} \approx \left( 
      \cos y_k -0 \right)  \Delta y.\] Summing the areas of all the slices and
    taking the limit as the number of slices $\to \infty$, the area of the
    region is 
    \begin{align*}
      \int_0^{\pi/2} \left( \cos y - 0 \right) dy &=
      \left. \sin y \right|_0^{\pi/2} \\
      &= \left( 
    \sin(\tfrac{\pi}{2}) - \sin(0) \right) \\
      &= \boxed{1-0 = 1}
    \end{align*}
  \end{solution}


\item

\begin{grading}
    Looking for mentions of slicing horizontally potentially reducing the number of integrals needed or transforming the integral into something easier to compute. They also should cite specific problems in their responses.
\end{grading}

  \begin{problem}[\vfill]
    \textbf{Reflect Back.}  Looking back at the problems you did on this
    assignment, in what situations can it be more useful or convenient to
    slice an area horizontally rather than vertically?
  \end{problem}  

  \begin{solution}
    There are regions like that in {{\#2}}, where if you
    slice vertically, there are multiple ``types'' of slices but, if you
    slice horizontally, there's just one ``type'' of slice.  In these
    situations, slicing horizontally allows you to write just one integral
    rather than a sum of multiple integrals.

    There are also situations like that in {{\#4}}, in
    which it's easy to write an integral by slicing vertically, but the
    integral is hard to evaluate.  Slicing horizontally allows you to
    write a different integral, which may be easier to evaluate. 
  \end{solution}

\end{enumerate}

\spaceonly{\psreflection{
Optional reading for next class is \S 27.1 in Gottlieb.
}}

\end{document}